\section{One-dimensional Noetherian domains}
\label{am-ch9-one-dimensional}
\begin{proposition}[Primary factorization in dimension one]
\label{am-prop-9-1}
\dcref{9.1}
\source{atiyah-macdonald-commutative-algebra:page-0102}
If $A$ is a Noetherian domain of dimension one, every nonzero ideal has a
unique expression as a product of primary ideals with distinct radicals.
\end{proposition}
\begin{proof}
Take a minimal primary decomposition.  Every nonzero prime is maximal, so the
distinct radicals are pairwise coprime; intersections therefore equal products.
Conversely a product of primary ideals gives the same minimal decomposition,
and isolated components are unique by Corollary~\ref{am-cor-4-11}.
\uses{am-thm-7-13,am-prop-1-16,am-prop-1-10,am-cor-4-11}
\end{proof}

\section{Discrete valuation rings}
\label{am-ch9-dvr}
\begin{definition}[Discrete valuation]
\label{am-def-discrete-valuation}
\source{atiyah-macdonald-commutative-algebra:page-0103}
A discrete valuation on a field $K$ is a surjection $v:K^*\to\Z$ with
$v(xy)=v(x)+v(y)$ and $v(x+y)\ge\min(v(x),v(y))$.  Its valuation ring is
$\{0\}\cup\{x:v(x)\ge0\}$.  A discrete valuation ring (DVR) is the valuation
ring of a discrete valuation on its fraction field.
\end{definition}
\begin{remark}[DVR ideals]
\label{am-rem-dvr-ideals}
\source{atiyah-macdonald-commutative-algebra:page-0103}
In a DVR, equal-valued elements differ by a unit; every nonzero ideal is
$\mathfrak m^n=(\pi^n)$ for a uniformizer $\pi$; the maximal ideal is
$(\pi)$ and is the unique nonzero prime.  Hence a DVR is a Noetherian local
domain of dimension one.
\end{remark}
\begin{proposition}[Characterizations of a DVR]
\label{am-prop-9-2}
\dcref{9.2}
\source{atiyah-macdonald-commutative-algebra:page-0103}
For a Noetherian local domain $(A,\m)$ of dimension one, the following are
equivalent: $A$ is a DVR; $A$ is integrally closed; $\m$ is principal;
$\dim_{A/\m}\m/\m^2=1$; every nonzero ideal is a power of $\m$; and there is
$x\in A$ with every nonzero ideal $(x^n)$.
\end{proposition}
\begin{proof}
DVRs are integrally closed by Proposition~\ref{am-prop-5-18}.  If $A$ is
integrally closed, choose $a\in\m$ of minimal valuation and $b\in\m^{n-1}$
outside $(a)$; the element $a/b$ forces $\m=(a/b)$.  A principal maximal ideal
gives one-dimensional $\m/\m^2$ by Nakayama, and conversely this dimension
and Proposition~\ref{am-prop-2-8} make $\m$ principal.  Applying the principal
ideal argument to $A/\m^n$ shows every ideal is a power of $\m$.  A generator
of $\m$ then gives all powers, and assigning to each nonzero element the
exponent of its principal ideal defines a well-defined discrete valuation.
\uses{am-prop-5-18,am-prop-2-8,am-cor-8-2,am-prop-8-6}
\end{proof}

\section{Dedekind domains}
\label{am-ch9-dedekind}
\begin{theorem}[Characterizations of Dedekind domains]
\label{am-thm-9-3}
\dcref{9.3}
\source{atiyah-macdonald-commutative-algebra:page-0104}
For a Noetherian domain $A$ of dimension one, the following are equivalent:
$A$ is integrally closed; every primary ideal is a power of a prime; and every
localization $A_\p$ at a nonzero prime is a DVR.  A ring with these properties
is a Dedekind domain.
\end{theorem}
\begin{proof}
The equivalence of the first and third conditions is Proposition~\ref{am-prop-9-2}
and local integral closedness (Proposition~\ref{am-prop-5-13}).  The second is
equivalent by localization of primary ideals and powers (Proposition~\ref{am-prop-4-8}
and Proposition~\ref{am-prop-3-11}).
\uses{am-prop-9-2,am-prop-5-13,am-prop-4-8,am-prop-3-11}
\end{proof}
\begin{corollary}[Unique ideal factorization]
\label{am-cor-9-4}
\dcref{9.4}
\source{atiyah-macdonald-commutative-algebra:page-0104}
Every nonzero ideal in a Dedekind domain factors uniquely as a product of
nonzero prime ideals.
\uses{am-prop-9-1,am-thm-9-3}
\end{corollary}
\begin{proof}
Use Proposition~\ref{am-prop-9-1} and replace each primary factor by its prime
power using Theorem~\ref{am-thm-9-3}.
\end{proof}
\begin{theorem}[Ring of integers]
\label{am-thm-9-5}
\dcref{9.5}
\source{atiyah-macdonald-commutative-algebra:page-0105}
The integral closure of $\Z$ in a number field is a Dedekind domain.
\end{theorem}
\begin{proof}
The extension is finite separable, so Proposition~\ref{am-prop-5-17} bounds
the integral closure inside a finite $\Z$-module; it is therefore Noetherian.
It is integrally closed by Corollary~\ref{am-cor-5-5}.  A nonzero prime
contracts nontrivially to $\Z$ by incomparability and is maximal by
Corollary~\ref{am-cor-5-8}; hence the dimension is one.
\uses{am-prop-5-17,am-cor-5-5,am-cor-5-8,am-cor-5-9}
\end{proof}

\section{Fractional ideals}
\label{am-ch9-fractional}
\begin{definition}[Fractional and invertible ideals]
\label{am-def-fractional-ideal}
\source{atiyah-macdonald-commutative-algebra:page-0105}
For a domain $A$ with fraction field $K$, an $A$-submodule $M\subseteq K$ is a
fractional ideal if $xM\subseteq A$ for some $0\ne x\in A$.  It is invertible
if $MN=A$ for some submodule $N$; then $N=(A:M)=\{z:zM\subseteq A\}$ and
is unique.  Nonzero principal fractional ideals are invertible.
\end{definition}
\begin{proposition}[Invertibility is local]
\label{am-prop-9-6}
\dcref{9.6}
\source{atiyah-macdonald-commutative-algebra:page-0106}
For a fractional ideal $M$, invertibility is equivalent to finite generation
and invertibility of every $M_\p$, and equivalently finite generation and
invertibility of every $M_\m$.
\end{proposition}
\begin{proof}
If $MN=A$, localize and use ideal quotients (Corollary~\ref{am-cor-3-15}) to
obtain $M_\p(A_\p:M_\p)=A_\p$.  Conversely let
$I=M(A:M)\subseteq A$.  Local invertibility gives $I_\m=A_\m$ for all maximal
$\m$, hence $I=A$ by local detection (Proposition~\ref{am-prop-3-9}); so
$M$ is invertible.
\uses{am-cor-3-15,am-prop-3-9}
\end{proof}
\begin{proposition}[Local criterion for a DVR]
\label{am-prop-9-7}
\dcref{9.7}
\source{atiyah-macdonald-commutative-algebra:page-0106}
For a local domain $A$, $A$ is a DVR iff every nonzero fractional ideal is
invertible.
\end{proposition}
\begin{proof}
In a DVR every fractional ideal is a power of a uniformizer.  Conversely,
invertible integral ideals are finite; choose a maximal non-power ideal $\aideal$.
Then $\m^{-1}\aideal$ is a strictly larger proper ideal unless $\aideal=0$,
contradicting maximality and Nakayama.  Thus all ideals are powers of $\m$ and
Proposition~\ref{am-prop-9-2} applies.
\uses{am-prop-2-6,am-prop-9-2}
\end{proof}
\begin{theorem}[Fractional ideal characterization]
\label{am-thm-9-8}
\dcref{9.8}
\source{atiyah-macdonald-commutative-algebra:page-0106}
For an integral domain $A$, $A$ is Dedekind iff every nonzero fractional ideal
is invertible.
\end{theorem}
\begin{proof}
In a Dedekind domain, fractional ideals are finite and localize to fractional
ideals of DVRs; apply Proposition~\ref{am-prop-9-7} and Proposition~\ref{am-prop-9-6}.
Conversely, invertible integral ideals are finite, so $A$ is Noetherian.  For a
nonzero prime $\p$, every integral ideal of $A_\p$ is the localization of an
invertible ideal and is invertible; Proposition~\ref{am-prop-9-7} makes
$A_\p$ a DVR, hence Theorem~\ref{am-thm-9-3} applies.
\uses{am-prop-9-6,am-prop-9-7,am-thm-9-3}
\end{proof}
\begin{corollary}[Ideal group and class group]
\label{am-cor-9-9}
\dcref{9.9}
\source{atiyah-macdonald-commutative-algebra:page-0107}
The nonzero fractional ideals of a Dedekind domain form a group under product.
The principal ideals form a subgroup; the quotient is the ideal class group, and
$1\to A^*\to K^*\to I\to\operatorname{Cl}(A)\to1$ is exact.
\uses{am-thm-9-8}
\end{corollary}
\begin{proof}
The inverse of $M$ is $(A:M)$; uniqueness and the exact sequence follow from
the definitions of principal ideals and units.
\end{proof}
